% !TEX root = ../main.tex
% ============================================================================
%  Pacotes
% ============================================================================

% --- Língua e tipografia ---------------------------------------------------
\usepackage[T1]{fontenc}
\usepackage[utf8]{inputenc}
\usepackage[english, portuguese]{babel}   % a última é a língua principal
\usepackage{csquotes}
\usepackage[sfdefault]{plex-sans}       % IBM Plex Sans (texto)
\usepackage[scale=0.95]{plex-mono}      % IBM Plex Mono (código, emails)
% Alternativas (trocar a linha do plex-sans):
%   \usepackage[sfdefault]{FiraSans}        % mais técnico/compacto
%   \usepackage[default]{sourcesanspro}     % mais neutro
%   \usepackage{plex-serif}                 % texto com serifa, mais clássico
\usepackage{microtype}

% --- Página -----------------------------------------------------------------
\usepackage[a4paper, margin=2.5cm, headheight=15pt]{geometry}
\usepackage{setspace}
\usepackage{indentfirst}
\usepackage{fancyhdr}
\usepackage{titlesec}
\usepackage[titles]{tocloft}
\usepackage[nottoc, notbib]{tocbibind}

% --- Figuras e tabelas ------------------------------------------------------
\usepackage{xcolor}
\usepackage{tikz}               % capa
\usepackage{graphicx}
\graphicspath{{imagens/}}       % \includegraphics{topologia} procura em imagens/
\usepackage{float}              % permite [H]
\usepackage{caption}
\usepackage{subcaption}
\usepackage{booktabs}
\usepackage{tabularx}
\usepackage{longtable}
\usepackage{multirow}
\usepackage{makecell}
\usepackage{pdfpages}           % \includepdf para anexar PDFs (ex.: enunciado)

% --- Matemática, unidades, listas ------------------------------------------
\usepackage{amsmath, amssymb}
\usepackage{siunitx}
\sisetup{output-decimal-marker={,}}
\usepackage{enumitem}

% --- Código -----------------------------------------------------------------
\usepackage{listings}

% --- Referências ------------------------------------------------------------
% style=ieee (numérico, comum em engenharia). Para APA: style=apa
\usepackage[backend=biber, style=ieee]{biblatex}

% --- Ligações (hyperref perto do fim; glossaries depois dele) ---------------
\usepackage{hyperref}
\usepackage[nameinlink]{cleveref}
\usepackage[acronym, toc, nonumberlist, nopostdot, nomain]{glossaries}
\makenoidxglossaries            % não precisa de makeglossaries / perl
